\documentclass[11pt]{article}
\usepackage[a4paper,margin=1in]{geometry}
\usepackage{fontspec}
\usepackage[boldfont=false]{xeCJK}
\setCJKmainfont{FandolSong-Regular.otf}
\usepackage{amsmath}
\usepackage{amssymb}
\usepackage{array}
\usepackage{booktabs}
\usepackage{graphicx}
\usepackage{adjustbox}
\usepackage{enumitem}
\setlist[itemize]{topsep=2pt,partopsep=0pt,parsep=0pt,itemsep=2pt,leftmargin=1.4em}
\setlist[enumerate]{topsep=2pt,partopsep=0pt,parsep=0pt,itemsep=2pt,leftmargin=1.6em}
\usepackage{parskip}
\usepackage{fvextra}
\fvset{breaklines=true,breakanywhere=true,fontsize=\small}
\setlength{\parindent}{0pt}

\begin{document}

\begin{center}
  {\LARGE\bfseries People in business}\\[0.35em]
  {\large Igcse Business Studies}
\end{center}
\vspace{0.6em}

\section*{Why motivation matters}

\textbf{Motivation} 激励 is the drive that makes people want to work hard. A well-motivated \textbf{workforce} 员工队伍 — all the \textbf{employees} 员工 in a business — brings real benefits:

\begin{itemize}
\item higher \textbf{productivity} 生产率 — each worker produces more,

\item lower \textbf{absenteeism} 缺勤 — fewer workers stay away from work,

\item lower \textbf{labour turnover} 员工流动率 — fewer workers leave, so the business spends less time and money hiring and training new staff.

\end{itemize}

\subsection*{Motivation theories}

Three thinkers explain what makes workers work hard.

\begin{center}
\adjustbox{max width=\linewidth}{%
\begin{tabular}{ll}
\toprule
\textbf{Thinker} & \textbf{Main idea} \\
\midrule
\textbf{Taylor} & Workers are mainly motivated by pay. Pay them more for more work (piece rate) and they work harder. \\
\textbf{Maslow} & People have a \textbf{hierarchy of needs} 需求层次: five levels, from basic pay and safety up to respect and reaching your full potential. You must meet the lower needs before the higher ones can motivate. \\
\textbf{Herzberg} & Two groups of factors. \textbf{Hygiene factors} 保健因素 (pay, conditions, rules) do not motivate, but poor ones cause unhappiness. \textbf{Motivators} 激励因素 (achievement and responsibility) truly drive people to work harder. \\
\bottomrule
\end{tabular}}
\end{center}

\subsection*{Financial methods of motivation}

These reward workers with money.

\begin{center}
\adjustbox{max width=\linewidth}{%
\begin{tabular}{ll}
\toprule
\textbf{Method} & \textbf{What it means} \\
\midrule
\textbf{wage} 工资 & regular pay, often by the hour or week \\
\textbf{time rate} 计时工资 & pay for each hour worked \\
\textbf{piece rate} 计件工资 & pay for each item made \\
\textbf{salary} 薪金 & a fixed yearly amount, paid monthly \\
\textbf{commission} 佣金 & pay based on how much you sell \\
\textbf{bonus} 奖金 & an extra payment for good work \\
\textbf{profit sharing} 利润分享 & workers get a share of the company's profit \\
\textbf{fringe benefits} 额外福利 & extras like a company car, free meals or health care \\
\bottomrule
\end{tabular}}
\end{center}

\subsection*{Non-financial methods of motivation}

These motivate without extra money.

\begin{itemize}
\item \textbf{job rotation} 工作轮换 — workers switch between different tasks, so work is less boring.

\item \textbf{job enlargement} 工作扩大化 — adding more tasks at the same level.

\item \textbf{job enrichment} 工作丰富化 — giving more challenging tasks and more responsibility.

\item \textbf{autonomy} 自主权 — letting workers make some of their own decisions.

\item \textbf{teamworking} 团队合作 — working in small groups.

\item \textbf{training} 培训 and chances for \textbf{promotion} 晋升 — helping workers improve and move up.

\end{itemize}

\section*{Organisation and management}

\subsection*{Organisational charts}

An \textbf{organisational chart} 组织结构图 is a diagram that shows how a business is organised — who reports to whom. It shows several ideas:

\begin{itemize}
\item \textbf{hierarchy} 层级 — the levels of staff, from the top down.

\item \textbf{chain of command} 指挥链 — the line along which orders pass down from managers to workers.

\item \textbf{span of control} 管理幅度 — the number of \textbf{subordinates} 下属 (people directly below) that one manager controls.

\item \textbf{delayering} 扁平化 — removing one or more levels to cut costs and speed up messages.

\end{itemize}

A long chain of command (a "tall" structure) slows messages down. A wide span of control gives one manager many people to watch.

\subsection*{Functions of management}

The five main jobs (functions) of \textbf{management} 管理 are:

\begin{itemize}
\item \textbf{planning} 计划 — setting aims and deciding how to reach them,

\item \textbf{organising} 组织 — arranging people and resources,

\item \textbf{coordinating} 协调 — making different parts work together,

\item \textbf{commanding} 指挥 — giving instructions and guiding staff,

\item \textbf{controlling} 控制 — checking that targets are met.

\end{itemize}

\subsection*{Leadership styles}

A \textbf{leadership style} 领导风格 is the way a manager leads people.

\begin{center}
\adjustbox{max width=\linewidth}{%
\begin{tabular}{lll}
\toprule
\textbf{Style} & \textbf{How it works} & \textbf{Best when} \\
\midrule
\textbf{autocratic} 专制型 & the leader decides alone and gives orders & a crisis, or with new and untrained staff \\
\textbf{democratic} 民主型 & the leader asks staff for ideas before deciding & staff are skilled and want a say \\
\textbf{laissez-faire} 放任型 & the leader sets goals, then lets staff work freely & creative, expert teams \\
\bottomrule
\end{tabular}}
\end{center}

\subsection*{Trade unions}

A \textbf{trade union} 工会 is a group of workers who join together to protect their shared interests. A union can bargain for better pay and conditions, support a worker in a dispute, and push for safer work.

\section*{Recruitment, selection and training}

\subsection*{The recruitment process}

\textbf{Recruitment} 招聘 is finding and attracting people to apply for a job. \textbf{Selection} 选拔 is choosing the best person from those who apply. The process starts with three documents:

\begin{itemize}
\item \textbf{job analysis} 工作分析 — studying what the job involves,

\item \textbf{job description} 工作描述 — a list of the duties and tasks of the job,

\item \textbf{person specification} 人员规格 — the skills, qualities and qualifications the right person needs.

\end{itemize}

\subsection*{Internal and external recruitment}

\begin{itemize}
\item \textbf{Internal recruitment} 内部招聘 fills the job with someone already in the business. It is cheaper and quicker, and the person is already known, but it brings in no new ideas.

\item \textbf{External recruitment} 外部招聘 fills the job from outside. It brings new ideas and skills, but it costs more and takes longer.

\end{itemize}

Common selection methods include the \textbf{application form} 申请表, the \textbf{CV} 简历 (a short summary of a person's education and work history), the \textbf{interview} 面试, and skills \textbf{tests} 测试.

\subsection*{Training}

Training improves workers' skills. There are three main types.

\begin{itemize}
\item \textbf{induction training} 入职培训 — given to new workers, to learn about the business and their role.

\item \textbf{on-the-job training} 在职培训 — learning while doing the job, beside an experienced worker.

\item \textbf{off-the-job training} 脱产培训 — learning away from the workplace, for example at a college.

\end{itemize}

\begin{center}
\adjustbox{max width=\linewidth}{%
\begin{tabular}{lll}
\toprule
\textbf{Type} & \textbf{Benefits} & \textbf{Limitations} \\
\midrule
on-the-job & cheap; uses real tasks & the trainer stops their own work; bad habits can pass on \\
off-the-job & expert teaching; no work distractions & costs more; the skills may not fit the exact job \\
\bottomrule
\end{tabular}}
\end{center}

\subsection*{Reducing the workforce}

Sometimes a business must cut staff numbers.

\begin{itemize}
\item \textbf{Redundancy} 裁员 — the job is no longer needed, for example because a machine now does it. It is not the worker's fault.

\item \textbf{Dismissal} 解雇 — the worker is removed for a good reason, such as breaking rules or poor work.

\end{itemize}

Most countries have \textbf{employment laws} 劳动法 that protect workers. They cover a fair \textbf{contract of employment} 雇佣合同, safe conditions, fair pay, and protection from being sacked unfairly.

\section*{Communication}

\subsection*{Why communication matters}

Good \textbf{communication} 沟通 means information is sent, received and understood correctly. It helps workers know what to do, cuts mistakes, and makes people feel valued.

The communication process has five parts: a \textbf{sender} 发送者, a \textbf{message} 信息, a \textbf{medium} 媒介 (the way it is sent), a \textbf{receiver} 接收者, and \textbf{feedback} 反馈 (the receiver's reply, which shows the message was understood).

\subsection*{Types of communication}

\begin{itemize}
\item \textbf{One-way communication} 单向沟通 — information flows in one direction only, with no reply (for example, a public notice).

\item \textbf{Two-way communication} 双向沟通 — the receiver can reply and give feedback. It is usually better, because it checks understanding.

\item \textbf{Internal communication} 内部沟通 — between people inside the business.

\item \textbf{External communication} 外部沟通 — between the business and outside people, like \textbf{customers} 顾客 or \textbf{suppliers} 供应商.

\end{itemize}

\subsection*{Barriers to communication}

A \textbf{communication barrier} 沟通障碍 is anything that stops a message getting through clearly. Common barriers include:

\begin{itemize}
\item a message that is too long or unclear,

\item a poor medium, such as a weak phone line,

\item language differences between sender and receiver,

\item a receiver who is not paying attention.

\end{itemize}

To reduce barriers: keep messages short and clear, choose the right medium, and ask for feedback to check the message was understood.


\end{document}
